\documentclass[lettersize,journal]{IEEEtran}

% ---- Unicode sanitization (IMPORTANT) ----
\DeclareUnicodeCharacter{202A}{}
\DeclareUnicodeCharacter{202C}{}
\DeclareUnicodeCharacter{200E}{}
\DeclareUnicodeCharacter{200F}{}

% ---- Math & theorem support ----
\usepackage{amsmath,amssymb,amsfonts}
\usepackage{amsthm}
\usepackage{bm}

\newtheorem{theorem}{Theorem}
\newtheorem{lemma}{Lemma}
\newtheorem{definition}{Definition}
\newtheorem{game}{Game}
\usepackage{amssymb}

% ---- Algorithms ----
\usepackage{algorithm}
\usepackage{algpseudocode}

% ---- Figures & tables ----
\usepackage{graphicx}

% Plot placeholder: draws the figure once cmd/plot has produced the file,
% and a framed note until then, so the labels resolve and the document
% compiles before the measurement run has finished.
\newcommand{\plotfig}[1]{%
  \IfFileExists{#1}{%
    \includegraphics[width=\linewidth]{#1}%
  }{%
    \fbox{\parbox[c][0.5\linewidth][c]{\dimexpr\linewidth-2\fboxsep-2\fboxrule\relax}{%
      \centering\ttfamily\footnotesize plot pending: #1}}%
  }%
}
\usepackage{subcaption}
\usepackage{booktabs}
\usepackage{tabularx}
\usepackage{multirow}
\usepackage{array}
\usepackage{adjustbox}

\usepackage{array}
\usepackage{pifont}
\newcommand{\cmark}{\ding{51}} % checkmark
\newcommand{\xmark}{\ding{55}} % cross mark

% ---- Text & formatting ----
\usepackage{cite}
\usepackage{textcomp}
\usepackage{url}
\usepackage{verbatim}
\usepackage{stfloats}
\usepackage{enumitem}
\usepackage{setspace}
\usepackage{bm}
\usepackage{xcolor}
\usepackage{listings}
\usepackage{makecell}

% ---- References ----
\usepackage[capitalize]{cleveref}

\hyphenation{op-tical net-works semi-conduc-tor IEEE-Xplore}

% updated with editorial comments 8/9/2021

\begin{document}

\title{ZK-Redact: Efficient Zero-Knowledge Authorized Redaction and Provenance Auditing for Permissioned Blockchains}

%\IEEEpubid{0000--0000/00\$00.00~\copyright~2021 IEEE}
% Remember, if you use this you must call \IEEEpubidadjcol in the second
% column for its text to clear the IEEEpubid mark.

\author{
\IEEEauthorblockN{Rangsimann Sattayarom, Ratchanon Wongwitutai, Natthawat Tungsriworakan, \\
Tagrid Chongkolrattanapond, Somchart Fugkeaw\\}
\IEEEauthorblockA{Sirindhorn International Institute of Technology, Thammasat University, Thailand\\
6622771747@g.siit.tu.ac.th,
6622780268@g.siit.tu.ac.th,
6622772364@g.siit.tu.ac.th,
guyhd9119@gmail.com,
somchart@siit.tu.ac.th}
}

\maketitle

\begin{abstract}
Redactable permissioned blockchains allow legitimate modification of
committed transactions, but controlled mutability raises two challenges:
verifying redaction authorization efficiently, and preserving auditable
evidence once the original ledger state has changed. Existing approaches
variously disclose authorization information, verify each proof
independently, or require costly traversal of redaction history. This
paper proposes \textbf{ZK-Redact}, a privacy-preserving framework for
authorized redaction and provenance auditing in permissioned
blockchains. Requesters prove compliance with a registered redaction
policy in zero knowledge, without revealing the private attributes that
satisfy it. Authenticated requests are pre-filtered and their proofs
verified through a batch-parallel workflow before redaction execution,
reducing per-request overhead. Each successful redaction yields an
authenticated provenance record binding the target transaction, its
previous and updated states, the applicable policy, and the verified
authorization evidence. These records are maintained in an authenticated
history structure, letting auditors verify current transaction integrity
and the sequence of authorized redactions with compact cryptographic
proofs.
\end{abstract}


\begin{IEEEkeywords}
Redactable Blockchain, Zero-Knowledge Proof,
Verifiable Redaction, Redaction Provenance, Integrity Auditing,
Permissioned Blockchain.
\end{IEEEkeywords}

\section{Introduction}

Permissioned blockchains provide shared, tamper-evident records for
multi-organization applications, yet strict immutability conflicts with
practical needs such as correcting erroneous data, removing sensitive
content, and satisfying privacy regulations. Redactable blockchains
resolve this by enabling authorized modification while preserving ledger
consistency, commonly via chameleon hashes (CHs) \cite{ref1,ref2,ref9}.
Recent work extends CH-based redaction with decentralized authorization,
fine-grained access control, anonymity, dynamic updates, and
verifiability \cite{ref4,ref6,ref14,ref15,ref16,ref20}.

Supporting \emph{large-scale} privacy-preserving redaction nonetheless
remains challenging. First, \emph{privacy-preserving authorization
becomes a verification bottleneck}. Anonymous mechanisms protect
requester identity and authorization information
\cite{ref4,ref14,ref17}, but processing cryptographic evidence
independently for each of many concurrent requests inflates verification
overhead. Existing scalability-oriented schemes optimize blockchain
structures or redaction processing \cite{ref10,ref14,ref21}, leaving
authorization under massive concurrency insufficiently addressed.
Second, \emph{redaction execution is computationally expensive at
scale}. Each CH modification requires authorized collision computation
and ledger-state updates \cite{ref1,ref2,ref6,ref19}, while policy
enforcement, distributed authorization, and accountability add further
processing \cite{ref15,ref20}. Handling validated requests
individually repeats chaincode, transaction, and ledger operations that
could otherwise be amortized.
Third, \emph{verifying redaction history is inefficient}. Since a valid
CH redaction preserves blockchain linkage, the current ledger state does
not reveal how a transaction evolved. Existing traceability and
auditing mechanisms \cite{ref10,ref13,ref16,ref20} still require
locating and examining each historical record, whereas efficient
auditing needs authenticated provenance supporting direct retrieval and
compact verification of both individual redactions and the continuity of
a transaction's history.
We therefore propose \textbf{ZK-Redact}, a privacy-preserving framework
for scalable redaction and provenance auditing in permissioned
blockchains. State- and policy-bound zero-knowledge proofs (ZKPs)
authorize redactions without exposing private credentials or attributes,
while logical proof shards and intra-shard batching enable parallel
verification. Verified requests are batch-processed for CH-based
redaction to amortize common ledger overhead, and transaction-specific
hash-linked provenance is Merkle-authenticated and blockchain-anchored
for efficient, independently verifiable auditing.

The main contributions are as follows:

\begin{itemize}
\item \textbf{Sharded Privacy-Preserving ZK Authorization:}
We combine state- and policy-bound ZK authorization with logical proof
sharding for parallel verification and optional intra-shard batching,
reducing verification bottlenecks while preserving request-level
authorization and failure isolation.

\item \textbf{Optimized Batch Redaction Processing:}
We separate ZKP verification from execution and batch verified requests
for redaction. State revalidation prevents stale updates, while common
validation, chaincode, and ledger costs are amortized without aggregating
request-specific CH computations.

\item \textbf{Authenticated Provenance and Verifiable Auditing:}
Transaction-specific hash-linked histories, Merkle-authenticated entries,
and blockchain anchors enable direct provenance retrieval and verification
of authorization, history completeness, ordering, and state continuity
without exhaustive ledger traversal.

\end{itemize}
\section{Related Work}
\label{sec:related}

Redactable blockchains have evolved from CH-based rewriting
\cite{ref9} to controlled redaction in consortium and IIoT settings
\cite{ref2}. Subsequent schemes distribute trapdoor control
\cite{ref1,ref6,ref19}, enforce attribute- or policy-based
authorization \cite{ref15}, and support dynamic or
trust-aware control \cite{ref5,ref16}. These mechanisms strengthen
redaction security and access control, but largely focus on secure
modification rather than scaling the complete authorization--redaction
pipeline under high request concurrency.

Privacy-preserving authorization and accountability introduce further
processing overhead. Huang \emph{et al.}~\cite{ref14} support scalable
anonymous redaction, PriChain~\cite{ref4} provides privacy-preserving
fine-grained control, and Wang \emph{et al.}~\cite{ref17} combine
anonymity with accountability for IIoT. However, these approaches do not
distribute large concurrent authorization workloads through sharded,
batch ZKP verification. ZK-Redact addresses this limitation by
partitioning independent proofs across logical shards for parallel
processing, with intra-shard batching when supported.

Scalability has also been studied through efficient data structures and
redaction mechanisms. Wu \emph{et al.}~\cite{ref10} use an extended
Merkle tree for inserted-data redaction, while Fugkeaw \emph{et
al.}~\cite{ref21} target large-scale GDPR-compliant blockchain e-KYC.
Other designs improve dynamic or distributed redaction
\cite{ref1,ref5,ref6,ref19}. Nevertheless, authorization, CH
adaptation, and ledger processing are commonly performed on a
per-request basis, repeatedly incurring shared processing costs.
ZK-Redact instead separates proof verification from redaction execution
and batch-processes verified requests to amortize common chaincode,
validation, and ledger operations.

Verifiability and traceability are addressed by VRBC~\cite{ref13},
which supports verifiable redaction and integrity auditing, and by
schemes providing trust-based traceability~\cite{ref16} or lightweight
storage and permission supervision~\cite{ref20}.
Authenticated structures such as
EMT~\cite{ref10} further improve verification of modified data.
However, efficiently retrieving and verifying the ordered redaction
history of a specific transaction remains insufficiently addressed,
particularly when authorization evidence, history completeness, and
state continuity must be verified together. 
\section{The Proposed ZK-Redact Framework}
\label{sec:scheme}

\subsection{System Model}
\label{sec:system}

ZK-Redact operates over a permissioned blockchain shared among multiple organizations, as showned in Fig.~\ref{fig:system_model}.

\begin{figure}[t]
\centering
\includegraphics[width=0.6\linewidth]{zk_redact_system_model.PNG}
\caption{System model of the proposed ZK-Redact framework}
\label{fig:system_model}
\end{figure}

\begin{enumerate}

\item \textbf{Requester (RQ):} a registered user or organization that submits transaction-redaction requests and generates a ZKP demonstrating compliance with the applicable policy without revealing private authorization information.

\item \textbf{Redaction Gateway (RG):} serves as the admission-control layer. It authenticates and pre-validates requests, filters stale or duplicate submissions, and removes unnecessary requester-identifying information before forwarding accepted requests for proof verification.

\item \textbf{Proof Verification Layer (PVL):} provides scalable ZKP authorization verification by distributing proofs across logical shards for parallel processing, with intra-shard batch verification when supported. Only verified requests proceed to redaction.

\item \textbf{Permissioned Blockchain (PB):} maintains transactions, redaction policies, and authenticated states. It executes verified CH-based redactions and batch-processes validated requests to reduce common chaincode and ledger overhead.

\item \textbf{Provenance Audit Index (PAI):} maintains authenticated transaction-specific redaction histories by linking successive states and anchoring Merkle-based commitments on the PB for efficient provenance retrieval and verification.

\item \textbf{Auditor (AU):} retrieves and independently verifies a target transaction's provenance, including redaction authorization, history integrity and completeness, and state continuity, without accessing private ZKP witnesses or unrelated blockchain records.

\end{enumerate}


\subsection{System Process}
\label{sec:system-process}

ZK-Redact operates through six phases as follows.
\medskip
\subsubsection*{\textbf{Phase 1: System Setup}}
\label{sec:phase1}

This phase initializes the cryptographic, authorization, blockchain,
proof-processing, and provenance components of ZK-Redact.

\smallskip
\noindent\textbf{Step 1: System and Entity Initialization}

The system administrator initializes
\begin{equation}
PP=(H,\mathsf{Sig},\mathsf{CH},\mathsf{ZK},
N,B,\Delta,B_R,\Delta_R)
\label{eq:pp}
\end{equation}
where $H$ is a collision-resistant hash function; $\mathsf{Sig}$,
$\mathsf{CH}$, and $\mathsf{ZK}$ denote the signature, chameleon-hash,
and zero-knowledge proof schemes; $N$ is the number of proof shards;
and $(B,\Delta)$ and $(B_R,\Delta_R)$ are the size and waiting bounds
for proof-verification and redaction batches.

Each requester $U_u$ generates and registers
\begin{equation}
(pk_u,sk_u)\leftarrow\mathsf{Sig.KeyGen}(1^\lambda),
\qquad
E_u=(ID_u,pk_u)
\label{eq:user-keygen}
\end{equation}
retaining $sk_u$ for request authentication. Similarly, auditor $AU_a$
generates $(pk_a,sk_a)\leftarrow\mathsf{Sig.KeyGen}(1^\lambda)$ and
registers
\begin{equation}
E_a^{\mathrm{AU}}
=
(AID_a,pk_a,\Pi_a^{\mathrm{audit}})
\label{eq:auditor-registration}
\end{equation}
where $\Pi_a^{\mathrm{audit}}$ defines its authorized scope for
provenance queries in Phase~6.

\smallskip
\noindent\textbf{Step 2: Redaction Policy Registration}

For each protected transaction class, the policy administrator defines
\begin{equation}
P_q=
(PID_q,\mathcal{A}_q,\mathcal{C}_q,
\mathcal{W}_q,\mathcal{V}_q)
\label{eq:redaction-policy}
\end{equation}
where $\mathcal{A}_q$, $\mathcal{C}_q$, $\mathcal{W}_q$, and
$\mathcal{V}_q$ specify the authorization predicates, additional
conditions, permitted redaction scope, and validity period. Its
commitment
\begin{equation}
C_{P_q}
=
H\!\left(
PID_q\parallel\mathcal{A}_q\parallel\mathcal{C}_q
\parallel\mathcal{W}_q\parallel\mathcal{V}_q
\right)
\label{eq:policy-commitment}
\end{equation}
is recorded on the PB and subsequently bound to the public ZKP
statement.

\smallskip
\noindent\textbf{Step 3: ZKP and Chameleon-Hash Setup}

The system generates
\begin{equation}
(pp_{\mathrm{zk}},vk_{\mathrm{zk}})
\leftarrow
\mathsf{ZK.Setup}(1^\lambda,\mathcal{R}_{\mathrm{red}})
\label{eq:zk-setup}
\end{equation}
where $\mathcal{R}_{\mathrm{red}}$ defines policy-compliant redaction
authorization; $pp_{\mathrm{zk}}$ is available to requesters and
$vk_{\mathrm{zk}}$ to the PVL. It also generates
\begin{equation}
(pk_{\mathrm{ch}},tk_{\mathrm{ch}})
\leftarrow
\mathsf{CH.KeyGen}(1^\lambda)
\label{eq:ch-setup}
\end{equation}
where $tk_{\mathrm{ch}}$ is retained by the trusted redaction component
for authorized CH adaptation.

For each redactable transaction, the PB initializes content
$m_i^{(0)}$ with randomness $r_i^{(0)}$ and commits
\begin{equation}
h_i^{\mathrm{ch}}
=
\mathsf{CH.Hash}
(pk_{\mathrm{ch}},m_i^{(0)},r_i^{(0)})
\label{eq:initial-ch-commitment}
\end{equation}
as its redactable ledger digest. Subsequent authorized redactions preserve
$h_i^{\mathrm{ch}}$ through CH adaptation in Phase~4.

\smallskip
\noindent\textbf{Step 4: Blockchain and Proof-Shard Initialization}

The PB deploys chaincode for entity/policy registration, audit
authorization, transaction updates, redaction results, and provenance
commitments. The PVL initializes
\begin{equation}
\mathbb{S}
=
\{\mathcal{S}_1,\ldots,\mathcal{S}_N\},
\label{eq:proof-shards}
\end{equation}
where each $\mathcal{S}_j$ maintains an independent proof queue.
Runtime shard assignment and batch formation are performed in Phase~3.

\smallskip
\noindent\textbf{Step 5: Provenance Audit Index Initialization}

For each transaction $TID_i$ with authenticated initial state
$T_i^{(0)}$, the PAI initializes
\begin{equation}
v_i=0,\qquad
\mathcal{H}_i^{(0)}=\emptyset,\qquad
d_i^{(0)}=H(T_i^{(0)})
\label{eq:initial-history}
\end{equation}
and computes
\begin{equation}
c_i^{(0)}
=
H(TID_i\parallel0\parallel d_i^{(0)})
\end{equation}
\begin{equation}
A_i^{(0)}
=
H(TID_i\parallel0\parallel c_i^{(0)})
\label{eq:initial-provenance}
\end{equation}
The initial PAI root is
\begin{equation}
R_{\mathrm{PAI}}^{(0)}
=
\mathsf{MerkleRoot}
\left(
\{A_i^{(0)}\}_{i\in\mathcal{I}}
\right)
\label{eq:initial-pai-root}
\end{equation}
where $\mathcal{I}$ is the set of transactions indexed by the PAI.
The PB anchors $R_{\mathrm{PAI}}^{(0)}$ as the trusted initial
provenance state.
\medskip

\subsubsection*{\textbf{Phase 2: Redaction Request and Privacy-Preserving Authorization}}
\label{sec:phase2}

This phase authenticates redaction requests, proves policy compliance
without revealing private attributes, and filters invalid, stale, or
duplicate requests before ZKP verification.

\smallskip
\noindent\textbf{Step 1: Redaction Request Construction}

Requests are indexed by $\rho$, with $i=\mathsf{tx}(\rho)$ denoting the
target transaction. For $TID_i$ at version $v_i$, requester $U_u$ constructs
\begin{equation}
R_\rho=
(TID_i,v_i,Loc_\rho,M_\rho,PID_\rho,t_\rho,n_\rho)
\label{eq:redaction-request}
\end{equation}
where $(Loc_\rho,M_\rho)$ specifies the target and modification,
$PID_\rho$ identifies policy $P_q$, and $(t_\rho,n_\rho)$ provides
freshness.

\smallskip
\noindent\textbf{Step 2: Policy-Bound ZK Authorization and Authentication}

Using the current policy version $v_{P_q}$ and commitment $C_{P_q}$,
the requester forms
\begin{equation}
x_\rho
=(TID_i,v_i,Loc_\rho,M_\rho,PID_\rho,v_{P_q},C_{P_q},t_\rho,n_\rho)
\end{equation}
\begin{equation}
\eta_\rho=H(x_\rho)
\end{equation}
and private witness
\begin{equation}
w_\rho=(Cred_u,\mathsf{attr}_\rho,\zeta_\rho)
\label{eq:zk-witness}
\end{equation}
containing the requester's credential, private authorization attributes,
and auxiliary evidence. The relation
$\mathcal{R}_{\mathrm{red}}(x_\rho,w_\rho)=1$ holds iff $Cred_u$ is valid
and non-revoked, $\mathsf{attr}_\rho$ satisfies the policy committed by
$C_{P_q}$, and $(Loc_\rho,M_\rho)$ is permitted.

The requester generates
\begin{equation}
\pi_\rho\leftarrow
\mathsf{ZK.Prove}(pp_{\mathrm{zk}},x_\rho,w_\rho)
\label{eq:zk-proof}
\end{equation}
and authenticates the complete request context by
\begin{equation}
\sigma_\rho=
\mathsf{Sig.Sign}
\left(
sk_u,H(R_\rho\parallel v_{P_q}\parallel C_{P_q}\parallel\eta_\rho)
\right)
\label{eq:request-signature}
\end{equation}
It submits
\begin{equation}
Q_\rho=
(R_\rho,v_{P_q},C_{P_q},\eta_\rho,\sigma_\rho,\pi_\rho)
\label{eq:redaction-query}
\end{equation}
to the RG.

\smallskip
\noindent\textbf{Step 3: Gateway Pre-Validation and Deduplication}

The RG retrieves $pk_u$ and verifies
\begin{equation}
\mathsf{Sig.Verify}
\left(
pk_u,H(R_\rho\parallel v_{P_q}\parallel C_{P_q}\parallel\eta_\rho),
\sigma_\rho
\right)=1
\label{eq:signature-verification}
\end{equation}
It recomputes $\eta_\rho=H(x_\rho)$, validates $(t_\rho,n_\rho)$,
confirms the current transaction and policy versions, and verifies
$(PID_\rho,v_{P_q},C_{P_q})$. Invalid or stale requests are rejected.

For an accepted request, the RG computes
\begin{equation}
DID_\rho=
H(TID_i\parallel v_i\parallel Loc_\rho\parallel M_\rho
\parallel PID_\rho\parallel v_{P_q})
\label{eq:deduplication-id}
\end{equation}
and accepts it only if $DID_\rho\notin\mathcal{D}$, where $\mathcal{D}$
contains pending or processed identifiers. The accepted $DID_\rho$ is
then inserted into $\mathcal{D}$. Its version bindings prevent duplicate
processing while allowing new requests after the relevant state advances.

\smallskip
\noindent\textbf{Step 4: Identity-Minimized Downstream Submission}

The RG removes explicit requester identity and constructs
\begin{equation}
Q_\rho^{*}=(DID_\rho,x_\rho,\eta_\rho,\pi_\rho)
\label{eq:identity-minimized-query}
\end{equation}
It retains the $DID_\rho\!\rightarrow\!U_u$ mapping for accountability,
while downstream components receive only $Q_\rho^{*}$. Thus, requester
identity is hidden downstream, but not from the RG. The accepted request
proceeds to Phase~3 for proof-shard assignment and ZKP verification.
\medskip


\subsubsection*{\textbf{Phase 3: Sharded Batch ZKP Verification}}
\label{sec:phase3}

This phase distributes accepted ZK proofs across logical shards for parallel verification, with intra-shard batching when supported.

\smallskip
\noindent\textbf{Step 1: Proof-Shard Assignment}

For each accepted $Q_\rho^{*}$, the PVL assigns
\begin{equation}
j_\rho=
1+\big(\mathsf{Int}(H(DID_\rho))\bmod N\big),
\qquad
Q_\rho^{*}\rightarrow\mathcal{S}*{j*\rho}
\label{eq:phase3-shard}
\end{equation}
providing uniform distribution in expectation across the $N$ shards. Each $\mathcal{S}_j$ maintains an independent queue and executes concurrently with other active shards.

\smallskip
\noindent\textbf{Step 2: Shard-Level Batch Formation}

In verification round $s$, each active shard forms
\begin{equation}
\mathcal{B}_j^{(s)}
=
{(x_\rho,\pi_\rho,Q_\rho^{*})}_{\rho\in\mathcal{Q}_j^{(s)}},
\qquad
|\mathcal{B}_j^{(s)}|\leq B
\label{eq:proof-batch}
\end{equation}
where $\mathcal{Q}_j^{(s)}$ denotes the selected requests. The batch closes when
\begin{equation}
|\mathcal{B}_j^{(s)}|=B
\quad\lor\quad
\tau_j^{(s)}\geq\Delta
\label{eq:batch-close}
\end{equation}
where $\tau_j^{(s)}$ is the oldest request's waiting time, thereby bounding delay under sparse arrivals.

\smallskip
\smallskip
\noindent\textbf{Step 3: Parallel ZKP Verification and Verified Request Pool}

Active shards process their batches concurrently. Each proof is verified as
\begin{equation}
V_\rho=
\mathsf{ZK.Verify}(vk_{\mathrm{zk}},x_\rho,\pi_\rho)\in\{0,1\}.
\label{eq:individual-zk-verification}
\end{equation}
When supported, a shard may instead apply
$\mathsf{ZK.BatchVerify}$ to its batch; otherwise, proofs are verified
individually, so native batching is optional. If an aggregate batch fails,
it is recursively divided to isolate invalid proofs. Successfully verified
requests form
\begin{equation}
\mathcal{P}_{\mathrm{ver}}
=
\{Q_\rho^{*}\mid V_\rho=1\},
\label{eq:verified-pool}
\end{equation}
while invalid requests are rejected and their pending identifiers released.


\subsubsection*{\textbf{Phase 4: Batch Redaction and Authenticated Provenance Update}}

This phase batch-processes verified requests, executes authorized CH
redactions, and updates authenticated transaction-specific provenance.
CH adaptation remains request-specific, while common blockchain and
provenance operations are amortized across each batch.

\smallskip
\noindent\textbf{Step 1: Batch Formation and Revalidation}

Verified requests are grouped into
\begin{equation}
\mathcal{B}_{R}^{(e)}
=
\{Q_\rho^{*}\}_{\rho\in\mathcal{Q}_R^{(e)}},
\qquad
|\mathcal{B}_{R}^{(e)}|\le B_R,
\label{eq:redaction-batch}
\end{equation}
which closes at size $B_R$ or after the oldest request waits $\Delta_R$.
Immediately before execution, the PB checks
\begin{equation}
\mathsf{Fresh}_\rho=
[v_i=v_i^{\mathrm{cur}}]\land
[v_{P_q}=v_{P_q}^{\mathrm{cur}}],
\label{eq:redaction-revalidation}
\end{equation}
and forms
$\mathcal{Q}_F^{(e)}
=\{\rho\in\mathcal{Q}_R^{(e)}:\mathsf{Fresh}_\rho=1\}$.
Requests sharing a transaction are ordered; since $x_\rho$ binds $v_i$,
at most one remains fresh per transaction in a round. The fresh set is
committed by
\begin{equation}
R_B^{(e)}
=
\mathsf{MerkleRoot}
\left(
\{H(DID_\rho\parallel\eta_\rho)\}_{\rho\in\mathcal{Q}_F^{(e)}}
\right).
\label{eq:redaction-batch-root}
\end{equation}

\smallskip
\noindent\textbf{Step 2: Authorized CH Redaction}

For each $\rho\in\mathcal{Q}_F^{(e)}$, the trusted redaction component
modifies $m_i^{(v_i)}$ to $m_i^{(v_i+1)}$ and computes
\begin{equation}
r_i^{(v_i+1)}
\leftarrow
\mathsf{CH.Adapt}
(tk_{\mathrm{ch}},m_i^{(v_i)},r_i^{(v_i)},m_i^{(v_i+1)})
\label{eq:ch-adapt}
\end{equation}
such that
\begin{equation}
\mathsf{CH.Hash}(pk_{\mathrm{ch}},m_i^{(v_i)},r_i^{(v_i)})
=
\mathsf{CH.Hash}(pk_{\mathrm{ch}},m_i^{(v_i+1)},r_i^{(v_i+1)}).
\label{eq:ch-collision}
\end{equation}
Thus, $T_i^{(v_i)}$ advances to $T_i^{(v_i+1)}$ while its CH-based
ledger commitment, block Merkle root, and blockchain links remain
unchanged. Batching amortizes common validation, chaincode, commitment,
and ledger operations.

\smallskip
\noindent\textbf{Step 3: Provenance Record Generation}

For each successful request $\rho$, ZK-Redact computes
\begin{equation}
d_i^{\mathrm{old},(v_i+1)}=H(T_i^{(v_i)}),\qquad
d_i^{\mathrm{new},(v_i+1)}=H(T_i^{(v_i+1)})
\label{eq:state-digests}
\end{equation}
and binds its authorization evidence by
\begin{equation}
C_\rho^{\mathrm{auth}}
=
H(DID_\rho\parallel\eta_\rho\parallel H(\pi_\rho)).
\label{eq:authorization-evidence}
\end{equation}
It constructs
\begin{equation}
\begin{split}
PR_i^{(v_i+1)}
=\big(&TID_i,v_i,v_i+1,PID_\rho,v_{P_q},DID_\rho,e,\\
&d_i^{\mathrm{old},(v_i+1)},d_i^{\mathrm{new},(v_i+1)},
C_\rho^{\mathrm{auth}},R_B^{(e)},t_\rho'\big),
\end{split}
\label{eq:provenance-record}
\end{equation}
where $(x_\rho,\pi_\rho)$ is retained off-chain under $DID_\rho$ and
authenticated by $C_\rho^{\mathrm{auth}}$.

\smallskip
\noindent\textbf{Step 4: Batch-Result Commitment}

Let $\mathcal{Q}_{\mathrm{succ}}^{(e)}\subseteq\mathcal{Q}_F^{(e)}$
be the successful requests and
$\mathcal{T}^{(e)}
=\{\mathsf{tx}(\rho):\rho\in\mathcal{Q}_{\mathrm{succ}}^{(e)}\}$.
ZK-Redact computes
\begin{equation}
\begin{aligned}
R_{PR}^{(e)}
&=
\mathsf{MerkleRoot}
\left(
\{H(PR_i^{(v_i+1)})\}_{i\in\mathcal{T}^{(e)}}
\right),\\
C_B^{(e)}
&=
H(R_B^{(e)}\parallel e\parallel R_{PR}^{(e)}).
\end{aligned}
\label{eq:batch-redaction-commitment}
\end{equation}
The PB applies the successful state updates; stale or failed requests
are excluded.

\smallskip
\noindent\textbf{Step 5: Transaction-Specific Provenance Update}

For each $i\in\mathcal{T}^{(e)}$, the PAI appends
$PR_i^{(v_i+1)}$ to $\mathcal{H}_i$ and updates
\begin{equation}
\begin{aligned}
c_i^{(v_i+1)}
&=H\!\left(c_i^{(v_i)}\parallel H(PR_i^{(v_i+1)})\right),\\
A_i^{(v_i+1)}
&=H(TID_i\parallel(v_i+1)\parallel c_i^{(v_i+1)}).
\end{aligned}
\label{eq:pai-update}
\end{equation}
History continuity requires
\begin{equation}
d_i^{\mathrm{old},(1)}=d_i^{(0)},\qquad
d_i^{\mathrm{new},(k)}=d_i^{\mathrm{old},(k+1)},
\quad 1\le k<v_i^{(e)},
\label{eq:history-continuity}
\end{equation}
where $d_i^{(0)}$ is the authenticated initial-state digest from
Phase~1. Unmodified PAI entries remain unchanged.

\smallskip
\noindent\textbf{Step 6: PAI Authentication and Anchoring}

After round $e$, the PAI computes
\begin{equation}
R_{\mathrm{PAI}}^{(e)}
=
\mathsf{MerkleRoot}
\left(
\{A_i^{(v_i^{(e)})}\}_{i\in\mathcal I}
\right)
\label{eq:pai-root}
\end{equation}
and the PB records
\begin{equation}
AC^{(e)}
=
H\!\left(
e\parallel R_{\mathrm{PAI}}^{(e)}
\parallel R_{\mathrm{PAI}}^{(e-1)}
\parallel C_B^{(e)}
\right)
\label{eq:pai-anchor}
\end{equation}
with
\begin{equation}
AR^{(e)}
=
\big(
e,R_{\mathrm{PAI}}^{(e)},R_{\mathrm{PAI}}^{(e-1)},
R_B^{(e)},R_{PR}^{(e)},C_B^{(e)},AC^{(e)}
\big).
\label{eq:anchor-record}
\end{equation}
This binds the authenticated provenance state to its predecessor and
redaction batch; $R_{PR}^{(e)}$ supports subsequent verification of
transaction-specific provenance records.
\medskip


\subsubsection*{\textbf{Phase 5: Provenance Retrieval and Audit Verification}}
\label{sec:phase5}

This phase enables an authorized auditor to retrieve and independently
verify a transaction's redaction history without traversing unrelated
blockchain records.

\smallskip
\noindent\textbf{Step 1: Auditor Authentication and Authorization}

For target $TID_i$ and PAI state $e$, auditor $AU_a$ signs
\begin{equation}
\begin{aligned}
R_a^{\mathrm{audit}} &= (AID_a, TID_i, e, t_a, n_a), \\
\sigma_a &= \mathsf{Sig.Sign}\big(sk_a, H(R_a^{\mathrm{audit}})\big).
\end{aligned}
\label{eq:auditor-request}
\end{equation}
The PB verifies $\sigma_a$ and freshness and requires
\begin{equation}
\mathsf{AuditAuth}(AID_a, TID_i, e, \Pi_a^{\mathrm{audit}}) = 1.
\label{eq:auditor-authorization}
\end{equation}

\smallskip
\noindent\textbf{Step 2: Provenance and Audit-Evidence Retrieval}

Let $\bar v_i$ be the version of $TID_i$ committed at PAI state $e$.
The PAI returns
\begin{equation}
\mathcal{E}_i^{(e)}
=
\left(
\mathcal{H}_i^{(\bar v_i)},\bar v_i,c_i^{(\bar v_i)},
\mu_i^{(e)},\{\psi_{i,k}\}_{k=1}^{\bar v_i}
\right),
\label{eq:audit-evidence}
\end{equation}
where $\mu_i^{(e)}$ authenticates
$A_i^{(\bar v_i)}$ under $R_{\mathrm{PAI}}^{(e)}$, and
\begin{equation}
\psi_{i,k}
=
\mathsf{MerkleProof}
\left(H(PR_i^{(k)}),R_{PR}^{(e_k)}\right)
\label{eq:record-path}
\end{equation}
authenticates record $PR_i^{(k)}$ in its redaction round $e_k$.
Using each $DID_{i,k}$, the auditor retrieves $(x_{i,k},\pi_{i,k})$
and the corresponding anchored $AR^{(e_k)}$, together with target
anchor $AR^{(e)}$ and authenticated boundary state digests.

\smallskip
\noindent\textbf{Step 3: Provenance Integrity and Completeness Verification}

The auditor computes
\begin{equation}
\widehat A_i
=
H(TID_i\parallel\bar v_i\parallel c_i^{(\bar v_i)})
\label{eq:audit-entry}
\end{equation}
and verifies its membership under $R_{\mathrm{PAI}}^{(e)}$ using
$\mu_i^{(e)}$. For each $1\le k\le\bar v_i$, it verifies
\begin{equation}
\mathsf{MerkleVerify}
\left(H(PR_i^{(k)}),\psi_{i,k},R_{PR}^{(e_k)}\right)=1
\label{eq:audit-record-merkle}
\end{equation}
and confirms that $R_{PR}^{(e_k)}$ belongs to the corresponding
blockchain anchor $AR^{(e_k)}$.

Starting from the authenticated $c_i^{(0)}$ of Phase~1, it reconstructs
\begin{equation}
\widehat c_i^{(k)}
=
H\left(
\widehat c_i^{(k-1)}\parallel H(PR_i^{(k)})
\right),
\quad
1\le k\le\bar v_i
\label{eq:audit-chain}
\end{equation}
with $\widehat c_i^{(0)}=c_i^{(0)}$. Integrity and completeness require
\begin{equation}
|\mathcal H_i^{(\bar v_i)}|=\bar v_i,\qquad
\widehat c_i^{(\bar v_i)}=c_i^{(\bar v_i)}
\label{eq:audit-completeness}
\end{equation}

\smallskip
\smallskip
\noindent\textbf{Step 4: Authorization, Continuity, and Anchor Verification}

For each $PR_i^{(k)}$, the auditor verifies
\begin{equation}
\begin{aligned}
H(DID_{i,k}\parallel H(x_{i,k})\parallel H(\pi_{i,k}))
&=C_{i,k}^{\mathrm{auth}},\\
\mathsf{ZK.Verify}(vk_{\mathrm{zk}},x_{i,k},\pi_{i,k})&=1,
\end{aligned}
\label{eq:audit-authorization}
\end{equation}
and checks consistency of its transaction, policy, and version information.
History continuity requires
\begin{equation}
d_i^{\mathrm{old},(1)}=d_i^{(0)},\qquad
d_i^{\mathrm{new},(k)}=d_i^{\mathrm{old},(k+1)},
\quad 1\le k<\bar v_i,
\label{eq:audit-continuity}
\end{equation}
with consecutive versions from $0$ to $\bar v_i$ and terminal digest
$d_i^{\mathrm{new},(\bar v_i)}$ matching the authenticated state at $e$.

For each required anchor $r\in\{e,e_1,\ldots,e_{\bar v_i}\}$, the auditor
further verifies
\begin{equation}
\begin{aligned}
AC^{(r)}
&\stackrel{?}{=}
H(r\parallel R_{\mathrm{PAI}}^{(r)}
\parallel R_{\mathrm{PAI}}^{(r-1)}\parallel C_B^{(r)}),\\
C_B^{(r)}
&\stackrel{?}{=}
H(R_B^{(r)}\parallel r\parallel R_{PR}^{(r)}).
\end{aligned}
\label{eq:audit-anchors}
\end{equation}
The audit succeeds only if the provenance membership, authorization,
history integrity and continuity, and blockchain anchors are valid,
enabling transaction-specific verification without exhaustive ledger
traversal.


\subsection*{Security Discussion}
\label{sec:security-discussion}

ZK-Redact relies on the security of the underlying signature, ZKP,
chameleon-hash, and collision-resistant hash schemes. Unauthorized
redaction is prevented by signature verification and ZK proofs of policy
compliance, while version binding, freshness checks, and pre-execution
revalidation prevent replay and stale requests. Zero knowledge protects
private authorization attributes, and requester identity is removed
before downstream verification. Authorized CH adaptation preserves
blockchain consistency without exposing the trapdoor. Finally, provenance
records bind each redaction to its authorization and state transition,
while hash-linked histories, Merkle commitments, and blockchain anchors
make modification, deletion, or reordering detectable during auditing.
Assuming secure CH-trapdoor management and permissioned consensus,
ZK-Redact therefore provides privacy-preserving authorization, controlled
redaction, and tamper-evident provenance.
\section{Performance Evaluation}

\subsection{Experimental Setup}
\label{sec:setup}

ZK-Redact and three baselines~\cite{ref10,ref13,ref22} are implemented
in Go at 128-bit security using SHA-256, Groth16 over BLS12-381,
P-256 chameleon hashing and Schnorr signatures, and a 3072-bit RSA
group for universal accumulation~\cite{ref22}. The ledger runs
Hyperledger Fabric with four organizations (two peers each) and a
three-node Raft ordering service, using 100 transactions/block and a
1-s block timeout. Experiments run on a 32-vCPU, 64-GB Ubuntu 22.04
host, with each peer limited to 1 vCPU and 2 GB RAM. The dataset contains 100{,}000 transactions with a 256-byte immutable
core and 512-byte redactable payload, 100 requesters, and 10 redaction
policies of predicate depth three. Each scheme replays 5{,}000 requests
under conflict ratios $\{0,0.25,0.50\}$. Off-chain cryptographic and
on-chain ledger costs are measured separately; baseline~\cite{ref10}
is evaluated both in-process and on Fabric to isolate consensus
overhead.

\begin{figure}[!t]
    \centering
    \includegraphics[width=0.32\columnwidth]{Images/1.jpeg}
    \hfill
    \includegraphics[width=0.32\columnwidth]{Images/2.png}
    \hfill
    \includegraphics[width=0.32\columnwidth]{Images/3.png}
    \captionsetup{justification=centering}
    \caption{Experimental results: (a) Verification Throughput (b) Redaction Processing Overhead (c) Provenance Retrieval and Audit Verification Cost}
    \label{fig:combined_3_panel}
\end{figure}

\subsubsection*{\textbf{Experiment 1: Verification Throughput}}
\label{exp:verification-throughput}

This experiment evaluates verification scalability as request concurrency
increases from $1$ to $512$. We compare ZK-Redact with in-process EMT,
Fabric-hosted EMT, and an ablation without sharding and batching.

As shown in Fig.~\ref{fig:combined_3_panel}(a), ZK-Redact and in-process
EMT achieve comparable throughput at low concurrency, but diverge beyond
approximately $32$ concurrent requests. EMT approaches saturation, whereas
ZK-Redact continues scaling by distributing proofs across shards and
amortizing verification within batches. In contrast, the ablation remains
nearly flat, confirming that the scalability gain stems from the proposed
sharding and batching rather than the underlying proof primitive.
Fabric-hosted EMT performs substantially worse because each verification
incurs chaincode and consensus overhead. Overall, ZK-Redact provides little
advantage for isolated requests, but increasingly outperforms the baselines
as concurrency grows, demonstrating its benefit under high verification
load.
\subsubsection*{\textbf{Experiment 2: Redaction Processing Overhead}}
\label{exp:redaction-overhead}

This experiment evaluates the average processing overhead per redaction
request as the number of requests increases, comparing ZK-Redact with
the three baselines~\cite{ref10,ref13,ref22}. As shown in Fig.~\ref{fig:combined_3_panel}(b), ZK-Redact consistently
achieves the lowest per-request overhead and remains nearly stable as the
workload grows. Ref.~\cite{ref10} incurs moderately higher cost, whereas
Refs.~\cite{ref13,ref22} are substantially more expensive, with
Ref.~\cite{ref22} exhibiting the largest overhead. The advantage of
ZK-Redact results from batch-oriented redaction, which amortizes common
validation, chaincode, and commitment-update costs across multiple
requests. Thus, ZK-Redact reduces redaction cost while maintaining stable
processing overhead under increasing workload.

\subsubsection*{\textbf{Experiment 3: Provenance Retrieval and Audit Verification Cost}}
\label{exp:provenance-audit}

This experiment evaluates provenance retrieval and audit verification as
the ledger grows from $10^2$ to $10^3$ blocks, comparing ZK-Redact with
the three baselines~\cite{ref10,ref13,ref22}.As shown in Fig.~\ref{fig:combined_3_panel}(c), ZK-Redact maintains low,
nearly constant audit latency as ledger size increases and remains
competitive with Refs.~\cite{ref10,ref13}. In contrast,
Ref.~\cite{ref22} incurs orders-of-magnitude higher cost that increases
with ledger size. This difference stems from the PAI, which retrieves
transaction-specific provenance and verifies it against authenticated
commitments without traversing the ledger history. The results therefore
show that ZK-Redact provides scalable provenance auditing while adding
only modest overhead relative to lightweight baselines.

\section{Conclusion}

This paper presented ZK-Redact, a privacy-preserving framework for scalable redaction and provenance auditing in permissioned blockchains. The proposed framework integrates state- and policy-bound zero-knowledge authorization, logical proof sharding with intra-shard batching for parallel verification, batch-oriented redaction processing that amortizes chaincode and ledger overhead,
and a Merkle-authenticated, blockchain-anchored Provenance Audit Index for efficient auditing. Experimental results demonstrated sustained verification throughput under high request concurrency, reduced redaction-processing overhead through batching, and audit verification cost that remains independent of ledger length, compared with representative redactable-blockchain baselines. 

\begin{thebibliography}{99}

\bibitem{ref1}
M. Jia, J. Chen, K. He, R. Du, L. Zheng, M. Lai, D. Wang, and F. Liu,
``Redactable Blockchain from Decentralized Chameleon Hash Functions,''
\textit{IEEE Trans. Inf. Forensics Security}, vol. 17, pp. 2771--2783, 2022,
doi: 10.1109/TIFS.2022.3192716.

\bibitem{ref2}
K. Huang, X. Zhang, Y. Mu, X. Wang, G. Yang, X. Du, F. Rezaeibagha, Q. Xia, and M. Guizani,
``Building Redactable Consortium Blockchain for Industrial Internet-of-Things,''
\textit{IEEE Trans. Ind. Informat.}, vol. 15, no. 6, pp. 3670--3679, Jun. 2019,
doi: 10.1109/TII.2019.2901011.

\bibitem{ref4}
H. Guo, W. Gan, M. Zhao, C. Zhang, T. Wu, L. Zhu, and J. Xue,
``PriChain: Efficient Privacy-Preserving Fine-Grained Redactable Blockchains in Decentralized Settings,''
\textit{Chinese Journal of Electronics}, vol. 34, no. 1, pp. 82--97, 2025.

\bibitem{ref5}
D. Zhang, J. Le, X. Lei, T. Xiang, and X. Liao,
``Secure Redactable Blockchain With Dynamic Support,''
\textit{IEEE Trans. Dependable Secure Comput.}, vol. 21, no. 2, pp. 717--731, 2024.

\bibitem{ref6}
X. Wu, X. Du, Q. Yang, N. Wang, and W. Wang,
``Redactable consortium blockchain based on verifiable distributed chameleon hash functions,''
\textit{J. Parallel Distrib. Comput.}, vol. 183, Art. no. 104777, 2024.

\bibitem{ref9}
G. Ateniese, B. Magri, D. Venturi, and E. R. Andrade,
``Redactable Blockchain -- or -- Rewriting History in Bitcoin and Friends,''
in \textit{Proc. 2nd IEEE Eur. Symp. Security Privacy (EuroS\&P)},
Paris, France, pp. 111--126, 2017.

\bibitem{ref10}
Z. Wu, L. Wang, X. Zhang and X. Feng, "EMT: Extended Merkle Tree Structure for Inserted Data Redaction in Permissioned Blockchain," in IEEE Transactions on Network Science and Engineering, vol. 12, no. 4, pp. 3025-3038, July-Aug. 2025.

\bibitem{ref13}
G. Tian, J. Wei, M. Kutylowski, W. Susilo, X. Huang, and X. Chen,
``VRBC: A Verifiable Redactable Blockchain with Efficient Query and Integrity Auditing,''
\textit{IEEE Trans. Comput.}, vol. 72, no. 7, pp. 1928--1942, Jul. 2023.

\bibitem{ref14}
K. Huang, X. Zhang, Y. Mu, F. Rezaeibagha, and X. Du,
``Scalable and redactable blockchain with update and anonymity,''
\textit{Information Sciences}, vol. 546, pp. 25--41, 2021.

\bibitem{ref15}
Y. Dong, Y. Li, Y. Cheng, and D. Yu,
``Redactable consortium blockchain with access control: Leveraging chameleon hash and multi-authority attribute-based encryption,''
\textit{High-Confidence Computing}, vol. 4, no. 1, Art. no. 100168, 2024.

\bibitem{ref16}
Y. Zhang, Z. Ma, S. Luo, and P. Duan,
``Dynamic Trust-Based Redactable Blockchain Supporting Update and Traceability,''
\textit{IEEE Trans. Inf. Forensics Security}, vol. 19, pp. 821--834, 2024.

\bibitem{ref17}
F. Wang, R. Dong, J. Cui, Q. Zhang and H. Zhong, "Blockchain-Based Anonymous and Accountable Data Redaction for the Industrial Internet of Things," in IEEE Transactions on Cloud Computing, doi: 10.1109/TCC.2026.3725253.

\bibitem{ref19}
X. Huang, Y. Wang, Y. Ding, Q. Wu, C. Yang, and H. Liang,
``Dynamically redactable Blockchain based on decentralized Chameleon hash,''
\textit{Digital Communications and Networks}, vol. 11, no. 3, pp. 757--767, 2025.

\bibitem{ref20}
M. Miao, X. Yang, J. Wei, G. Tian, and W. Susilo,
``A Verifiable and Redactable Blockchain with Lightweight Storage and Permission Supervision,''
\textit{Information}, vol. 17, no. 2, Art. no. 176, 2026.

\bibitem{ref21}
S. Fugkeaw, S. Sungchai, S. Nakprame, and P. Sreekongpan,
``Enabling Secure and Scalable GDPR-Compliant Blockchain-Based e-KYC With Efficient Redaction,''
\textit{IEEE Access}, vol. 13, pp. 136834--136853, 2025.

\bibitem{ref22}
Y. Shen,et al.,
``Verifiable and Redactable Blockchains With Fully Editing Operations,''
\textit{IEEE Transactions on Information Forensics and Security}, vol. 18, pp. 3792--3805, 2023.

\end{thebibliography}

\end{document}